\chapter{Discussion and Limitations}
\label{ch:xdiscussion}

The chapters so far share a pattern: each new capability is first a
double-precision tool tested against closed forms, then a live rendering
path compared with that tool. The comparisons are what make the images
defensible, and they have found real defects. Three were in the renderer:
the double-counted colour, the step budget and the missing constraint
projection. Two were in the orbit fitting. The forecast Jacobian needs
finite-difference steps far smaller than the thesis used, and the thesis'
direct no-hair fit cannot move $\epsq$ off its starting value
(Chapter~\ref{ch:xforecast}). Several limitations remain and are stated
where they arise.

\begin{itemize}
\item \emph{The thin disk is hypothetical.} At \sgra's measured accretion
rate a thin, radiatively efficient disk would be $\XLumFactorLow$--$\XLumFactorHigh$
times too luminous (\S\ref{sec:thindisk}). The renders that show one are
labelled as such; the physical flow is the subject of
Chapter~\ref{ch:xriaf}.
\item \emph{The sky is Earth's.} The Milky Way panorama is the sky seen from
the Solar System, used as a backdrop at infinity, and its colour shift treats
each pixel as a 6500\,K blackbody (\S\ref{sec:sky}).
\item \emph{The hot spot is the simplest model.} A rigid Gaussian on a
circular geodesic, optically thin, with a fixed spectral index. It
reproduces the GRAVITY periods and sense of rotation but not the size of the
observed centroid motion, exactly the tension GRAVITY reported. Shear, cooling,
out-of-plane motion and polarisation are not modelled
(Chapter~\ref{ch:xflares}).
\item \emph{Display processing is not physics.} Bloom, tone curves, white
balance and false colour are confined to the display chain and are switched
off in every quantitative comparison (\S\ref{sec:display}).
\item \emph{The accretion flow is a semi-analytic RIAF, not a simulation.}
Its radial profiles are fixed power laws, its electrons are purely thermal,
and its one free parameter is fitted to a single flux density. The
comparison with the EHT ring is therefore a test of this family of models,
not of general relativity. In that test the family's reference temperature
fails: the ring is too small. The thesis' critical-curve metrology, which
needs no emission model, remains the result that tests the spacetime.
\item \emph{Polarisation without Faraday effects.} Rotation and conversion in
the flow are omitted, and the intrinsic polarisation fraction is a fixed
assumption. Only quantities independent of that fraction --- the phase of
$\beta_2$ and $|\beta_2|/|\beta_1|$ --- are compared with the EHT. Interstellar
scattering toward \sgra{} is not modelled; the EHT notes that it barely
affects $\beta_2$ \citep{EHT2024pol}.
\end{itemize}

\begin{itemize}
\item \emph{The forecasts are idealised.} They use noise-free data, sampled
uniformly with no seasonal gaps, and include no reference-frame systematics,
source confusion or second star. They are local (Fisher) estimates, checked
against full profile-likelihood fits to within $\XFcProfAgree$ per cent. The
extended mass is a single Plummer profile at GRAVITY's scale radius, and the
spin axis lies along the line of sight. Both choices change the numbers,
though not the conclusion that an extended mass must be fitted with any
future no-hair or spin test (Chapter~\ref{ch:xforecast}).
\item \emph{The wormhole is not \sgra.} The Dneg metric is a test of the
renderer's generality. The exotic matter it would need is not modelled, its
other universe's sky is an illustrative recolouring, and nothing about the
existence of such objects is implied (Appendix~\ref{app:xwormhole}).
\item \emph{The fall is classical and stops at $r_-$.} It is one geodesic in
vacuum Kerr: no accretion flow, no quantum effects, and no interior beyond
the inner horizon, where mass inflation makes the classical solution
unreliable (Chapter~\ref{ch:xinfall}).
\end{itemize}
